\documentclass{article}
\usepackage[T1]{fontenc}
\usepackage{lmodern}
\usepackage{graphicx}
\usepackage{amsmath,amssymb}
\usepackage[a4paper,margin=2cm]{geometry}
\usepackage[absolute,overlay]{textpos}
\setlength{\TPHorizModule}{1mm}
\setlength{\TPVertModule}{1mm}
\usepackage[indonesian]{babel}
\usepackage{hyperref}
\setlength{\emergencystretch}{2em}
% unit: Halaman 30

\begin{document}

% === Halaman 30 (llm) ===

\section*{Penjumlahan Dua Titik di dalam EC pada GF(p)}

Misalkan $P(x_p, y_p)$ dan $Q(x_q, y_q)$. \\
Penjumlahan: $P + Q = R$

\subsection*{Koordinat Titik R:}
\[
\begin{aligned}
x_r &= m^2 - x_p - x_q \pmod{p} \\
y_r &= m(x_p - x_r) - y_p \pmod{p}
\end{aligned}
\]

$m$ adalah gradien:
\[
 m = \frac{y_p - y_q}{x_p - x_q} \pmod{p}
\]

\end{document}
